\documentclass[11pt,a4paper,sans]{moderncv}
\pagestyle{empty}
\moderncvstyle{classic} % styles: casual, classic, banking, oldstyle, fancy
\moderncvcolor{blue}    % colors: blue, orange, green, red, purple, grey, black

\usepackage{enumitem}
\usepackage{multicol}
\setlength{\columnsep}{3em}
\usepackage[utf8]{inputenc}
\usepackage[T1]{fontenc}
\usepackage[english]{babel}
\usepackage[a4paper, top=1.5cm, bottom=1.5cm, left=1.5cm, right=1.5cm]{geometry}

\name{Raphaël}{Jontef}
\photo[64pt][0.4pt]{../../photo.jpg}
\phone[mobile]{+33 7 68 14 51 09}
\email{raphaeljontef@hotmail.com}

\begin{document}
\makecvtitle

\section{Objectif}
Postuler à un stage de vulgarisation scientifique à destination d'élèves de seconde, plutôt dans les thématiques de l'intelligence artifielle ou de l'algorithmique.

\section{Formation}
\cventry{Paris 20\textsuperscript{e} arr.}{sept. 2020 -- juil. 2023}{Lycée Maurice Ravel}{}{}{Spécialités NSI et Mathématiques, mention "bien" au baccalauréat général et technologique.}
\vspace{1em}
\cventry{Paris 8\textsuperscript{e} arr.}{sept. 2023 -- juil. 2025}{Lycée Fénelon Sainte-Marie}{CPGE}{}{MP2I puis MPI (Mathématiques, Physique et Informatique)}
\vspace{1em}
\cventry{Bordeaux, Talence}{sept. 2025 -- 2028}{Enseirb-Matmeca}{école d'ingénieurs}{}{Filière Informatique}.

\section{Expériences professionnelles et projets}
\vspace{0.5em}
\cvitem{Enseignement}{Tutorat présenciel en sciences pour des élèves de collège et lycée, chez \href{https://www.alveusclub.com/}{Alveus}. Une cinquantaine d'heures de cours dispensées de novembre 2025 à aujourd'hui, principalement en physique et en mathématiques.}
\cvitem{Web}{Création de sites en \texttt{HTML}, \texttt{CSS}, \texttt{Javascript} dans un cadre aussi bien personnel que scolaire.}
\vspace{0.5em}
\cvitem{Docs}{Rédaction en \LaTeX\, de multiples documents (rapport et présentation de TIPE en CPGE, ce CV, prises de notes, rapports de projets en école, \textit{etc.}).}
\vspace{0.5em}
\cvitem{Autre}{Travail d'Initiative Personnelle Encadrée (TIPE) en CPGE, sur le thème de la reconnaissance de caractère par ordinateur pour retrouver du code source \LaTeX\,à partir d'une image.}

\section{Compétences informatiques}
\begin{itemize}[left=0pt, label=\textbullet]
\begin{multicols}{2}
    \item \texttt{Python}, \texttt{C} et \texttt{OCaml};
    \item Création de sites web~: \texttt{HTML}, \texttt{CSS} et \texttt{Javascript}, notamment avec \texttt{Vue.js};
    \item Écriture de documents en \LaTeX;
    \item utilisation de \texttt{git}.
\end{multicols}
\end{itemize}

\noindent
\begin{minipage}[t]{0.48\textwidth}  % gauche
\section{Compétences humaines}
\begin{itemize}[left=0pt, label=\textbullet]
    \item Communication bonne, notamment à l'écrit;
    \item Pédagogie (expérience de tutorat);
    \item Sens affûté du détail;
    \item Importante volonté de bien faire.
\end{itemize}
\end{minipage}
\hfill
\begin{minipage}[t]{0.48\textwidth}  % droite
\section{Passe-temps}
\begin{itemize}[left=0pt, label=\textbullet]
    \item Cuisine et patisserie;
    \item Cyclisme, aussi bien par besoin que par loisir;
    \item Algorithmique numérique et théorie algébrique associée;
\end{itemize}
\end{minipage}

\end{document}
